\scanpage{58}
% archive sheet number: 57
Soit $\underline{E}$ faisceau loc. libre de rang $r$. Alors \deleted{$\unsure{\operatorname{GP}(\underline{E})}$ \unsure{opère sur}} $P(S^n(\underline{E}))$, \inserted{et $P(S^n(\underline{E})^{\vee})=P(S^n(\unsure{\check{\underline{E}}|\underline{E}^{\vee}}))$} \unsure{ne dépend que de} $P(\underline{E})$, \deleted{\unsure{de façon précise}} et $\operatorname{GP}(\underline{E})$ opère sur $P(S^n(\underline{E}))$ \deleted{\unsure{de cette manière}} \inserted{$P({\unsure{\widetilde{S}|S}}^{n}(\underline{E})^{\vee})$} [\unsure{hom.}
\[
\operatorname{GP}(\underline{E})\longrightarrow\operatorname{GP}(S^n_{\underline{O}_{Y}}(\underline{E}))
\]
] Faisant $\underline{E}=\underline{O}_{Y}^{r+1}$, \unsure{on trouve que} $\operatorname{GP}(r)_{Y}$ opère dans $P(S^n(O_{Y}^{r+1})^{\vee})$, \unsure{d'où} \inserted{et par suite pour tout schéma de Br.-Sév. $X$} fibré principal \unsure{hom.} sous $\operatorname{GP}(r)_{Y}$, \deleted{\unsure{un fibré}\ \unsure{associé}} un fibré \deleted{\unsure{proj.}} \inserted{associé}, $\unsure{S|\mathscr{S}}$-isomorphe à $P(S^n(\underline{O}_{Y}^{r+1}))$, qui est donc un schéma de Brauer-Sévéri. Il est loc. triv. (trivial) si $X$ \deleted{l'est, \unsure{on le désigne}} \inserted{On l'appelle le} \uline{schéma des $n$-formes de} $X$, on le \unsure{dénote} par $\Phi^n(X)$. Compatibilité aux images inverses.
\margin{\unsure{n?}\ldots}
\margin{$\unsure{\operatorname{GP}(r)_{Y}}\longrightarrow\unsure{\operatorname{GP}}(S^n_{\underline{O}_{Y}}(O_{Y}^{r+1})^{\vee})\ \deleted{\unsure{\longrightarrow}}$}

Prenons d'abord $X=P(\underline{E})$. Alors
\[
\Phi^n(X)\overset{\text{can.}}{\simeq}P(\underline{S}^{n}_{\underline{O}_{X}}(\overset{\deleted{\unsure{\vee|\wedge}}}{\underline{E}})^{\vee}).
\]
[\unsure{démonstration} : \unsure{c'est} un foncteur, compatible aux images inverses\ldots, qui définit une \unsure{représentation} et par suite une notion de fibré associé; \unsure{dans} cette \deleted{\unsure{fibré}} notion d'associé \unsure{redonne} le foncteur].

\deleted{\unsure{Sections de} $\unsure{\Phi}$} Alors une section de $\Phi^{(n)}(X)$ est donc définie par \deleted{\illegible{1}} la donnée d'un \deleted{\unsure{épi}}\inserted{homo.}\unsure{morphisme} de faisceaux de $\underline{S}^{n}(\underline{E})^{\vee}$ sur un \deleted{\unsure{loc. lib.}} faisceau inversible, \unsure{donc correspondant} à un hom ``ne s'annulant nulle part'' d'un fibré inversible dans $\underline{S}^{n}(\underline{E})$, i.e. une section de
\[
\underbrace{\deleted{\unsure{\mathscr{H}}}(\underline{S}^{n}(\underline{E}))^{*}}_{\inserted{\substack{\text{faisceau des sections de }\underline{S}^{n}(\underline{E})\\\text{qui ne ``s'annulent pas''}}}}/\underline{O}_{Y}^{*};
\]
On \unsure{aura}
